\section{Methodology: Gradient Decoupling and the Feedback Firewall}
\label{sec:method}

In this section, we dissect the internal gradient mechanics of recurrent
feature aggregation, prove the mathematical paralysis inherent in conventional
coupled gating, and formulate our proposed \textbf{Feedback Firewall}
and \textbf{Detached Soft-OR} architecture.


% -----------------------------------------------------------------------------
\subsection{The Classical MixPool: Architecture and The Hard-Gating Flaw}
The Feature Attention Network incorporates recurrent refinement across
four encoder stages ($e_1, e_2, e_3, e_4$) and four decoder stages
($d_1, d_2, d_3, d_4$). At each level, intermediate feature tensors
$x \in \mathbb{R}^{B \times C \times H \times W}$ and the recurrent
feedback mask
$m \in \mathbb{R}^{B \times 1 \times H_{\text{in}} \times W_{\text{in}}}$
are processed through a \texttt{MixPool} block.

Within \texttt{MixPool}, input feature $x$ is first routed to a
lightweight convolutional attention sub-network $\mathcal{F}_{\text{att}}$,
generating a continuous single-channel attention prior:
\begin{equation}
  \label{eq:fmask_def}
  \begin{split}
    \text{fmask} &= \sigma\!\Big(
      \text{Conv}_{1 \times 1}\!\Big(
        \text{ReLU}\!\big(
          \text{BN}\!\big(\text{Conv}_{3 \times 3}(x)\big)
        \big)
      \Big)
    \Big) \\
    &\quad \in [0, 1]^{B \times 1 \times H \times W}.
  \end{split}
\end{equation}
Simultaneously, the feedback mask $m$ is spatially downsampled via max
pooling to match the spatial resolution $(H, W)$ of the current layer,
yielding $m_{\text{fg}} \in [0, 1]^{B \times 1 \times H \times W}$.

In the original formulation~\cite{tomar2021fanet}, these two spatial
gates are fused into a single binary selection mask
$\text{keep}_{\text{hard}}$ via a hard thresholding operator:
\begin{equation}
  \label{eq:keep_hard}
  \text{keep}_{\text{hard}} =
    \max\!\left(\mathbb{I}(\text{fmask} > 0.5),\; m_{\text{fg}}\right).
\end{equation}

The gated features are then split across dual convolutional pathways
and recombined:
\begin{equation}
  x_1 = \text{Conv}_1\!\left(x \odot \text{keep}_{\text{hard}}\right),
  \qquad
  x_2 = \text{Conv}_2(x),
\end{equation}
\begin{equation}
  \text{MixPool}(x, m) =
    \left[x_1 \,\|\, x_2\right]
    \in \mathbb{R}^{B \times C \times H \times W}.
\end{equation}

\subsubsection{The Calculus of Failure: Zero-Gradient and Toxic Coupling}
Let $\mathcal{L}$ denote the scalar segmentation objective. Applying the
chain rule to backpropagate gradients from $x_1$ back to the parameters
$\theta_{\text{att}}$ of the attention generator $\mathcal{F}_{\text{att}}$:
\begin{align}
  \label{eq:chain_rule}
  \frac{\partial \mathcal{L}}{\partial \theta_{\text{att}}}
    &= \frac{\partial \mathcal{L}}{\partial x_1}
       \cdot \frac{\partial x_1}{\partial \text{keep}_{\text{hard}}}
       \cdot \frac{\partial \text{keep}_{\text{hard}}}{\partial \text{fmask}}
       \cdot \frac{\partial \text{fmask}}{\partial \theta_{\text{att}}}.
\end{align}

Because $\text{keep}_{\text{hard}}$ utilizes the indicator function
$\mathbb{I}(\text{fmask} > 0.5)$, the distributional derivative of the
Heaviside step function is the Dirac delta:
\begin{align}
  \frac{\partial}{\partial u}\mathbb{I}(u > 0.5)
    &= \delta(u - 0.5) = 0 \quad \forall\, u \neq 0.5.
\end{align}
Consequently, the third factor in Eq.~(\ref{eq:chain_rule}) evaluates to:
\begin{align}
  \frac{\partial \text{keep}_{\text{hard}}}{\partial \text{fmask}}
    &\equiv 0 \quad \text{almost everywhere.}
\end{align}
As a direct consequence, \textbf{zero gradient reaches the parameters
$\theta_{\text{att}}$ of the attention generator}. The network cannot
adapt its internal spatial feature selector in response to task-specific
segmentation errors.

Simultaneously, the recurrent branch receives backward gradients through
automatic differentiation:
\begin{align}
  \frac{\partial \mathcal{L}}{\partial m_{\text{fg}}}
    &= \frac{\partial \mathcal{L}}{\partial x_1}
       \cdot \frac{\partial x_1}{\partial \text{keep}_{\text{hard}}}
       \cdot \frac{\partial \text{keep}_{\text{hard}}}{\partial m_{\text{fg}}}
       \neq 0.
\end{align}
When the network produces an early false-positive prediction, this
erroneous mask is reinjected into the encoder. Because gradients flow
through $m_{\text{fg}}$ back into the recursive history without internal
attention guidance, the encoder parameters are updated to maximize
consistency with past errors rather than anatomical boundaries. This
causes the internal feature representations of healthy background mucosa
to collapse into lesion-like embeddings.

% -----------------------------------------------------------------------------
\subsection{The Monotonicity Trap of Mathematical Gating}
To eradicate the Feedback Trap, our initial design (Phase 7B) explored a \textbf{Detached Soft-OR Gating} mechanism:
\begin{equation}
  \text{keep}_{\text{soft\_or}} = 1.0 - (1.0 - a) \odot (1.0 - \text{detach}(m_{\text{fg}})),
\end{equation}
where $a$ is the internal spatial attention map and $m_{\text{fg}}$ is the recurrent feedback mask. While gradient detachment successfully severed the toxic recursive gradient loop, empirical temporal tracking (FPR evaluation across iterations $t=0$ to $t=3$) revealed a catastrophic flaw: the False Positive Rate exploded to $52.31\%$.

We diagnose this failure as the \textbf{Soft-OR Monotonicity Trap}. The Soft-OR gate $A \lor B$ is mathematically a strictly monotonically increasing operator with respect to the prior mask $B$. Given that the initialization mask ($t=0$, generated via Otsu thresholding) frequently yields massive false positives in boundary-ambiguous flat polyps, the Soft-OR gate ensures that $\text{keep}_{\text{soft\_or}} \ge m_{\text{fg}}$. 
Consequently, the network is physically paralyzed---it is mathematically incapable of subtracting, pruning, or suppressing the noisy false positives inherited from previous iterations.

% -----------------------------------------------------------------------------
\subsection{Proposed Solution: Phase 7C Learned Residual Decoupling}
To overcome the Monotonicity Trap, the gating mechanism must be fully differentiable, structurally decoupled from recursive gradient flow, and capable of both augmenting and \textit{pruning} the prior mask.

We propose \textbf{Learned Residual Decoupling}. We discard rigid Boolean mathematics and instead introduce a low-parameter dynamic $1\times 1$ convolutional spatial attention gate:
\begin{equation}
  \text{combined} = [a \;\|\; \text{detach}(m_{\text{fg}})],
\end{equation}
\begin{equation}
  G = \sigma\Big(\text{Conv}_{1 \times 1}(\text{combined})\Big),
\end{equation}
where $G \in (0, 1)^{B \times 1 \times H \times W}$ acts as a spatial modulating gate. The final feature selection tensor is computed via a residual addition:
\begin{equation}
  \label{eq:learned_residual}
  \text{keep}_{\text{proposed}} = (G \odot \text{detach}(m_{\text{fg}})) + a.
\end{equation}

\subsubsection{Theoretical Advantages}
Our Phase 7C architecture possesses three critical properties:
\begin{enumerate}
    \item \textbf{Non-Monotonicity (Pruning Capability):} Unlike the OR-gate, if the network detects a false positive in the recurrent prior $m_{\text{fg}}$, it can actively push the learned gate $G \to 0$ and $a \to 0$, forcing $\text{keep}_{\text{proposed}} \to 0$. This fully restores the network's ability to prune historical errors.
    \item \textbf{Gradient Preservation:} Because the $1 \times 1$ convolution is continuous and differentiable, dense gradients flow back into the feature mask $a$.
    \item \textbf{Structural Decoupling:} We strictly preserve $\text{detach}(m_{\text{fg}})$. The recurrent gradient loop remains severed, forcing the internal feature encoder $a$ (and the learned gate $G$) to take absolute responsibility for minimizing false positives.
\end{enumerate}

To ensure the network utilizes its new pruning capabilities, we combine this architecture with a \textbf{Curriculum Adaptive Tversky Loss}, placing heavy asymmetric penalties on False Positives ($\alpha=0.7$, $\beta=0.3$) once the network reaches basic stability.
